\BourbakiExerciseGroup

\begin{enumerate}
\def\labelenumi{\arabic{enumi})}
\tightlist
\item
  在拓扑群 G 上，右一致结构是唯一的与 G 的拓扑相容且具有一个一致邻域基本系（其中每个一致邻域都在所有平移下不变）的一致结构
\end{enumerate}

\[
(x, y) \rightarrow (x a, y a).
\]

\begin{enumerate}
\def\labelenumi{\arabic{enumi})}
\setcounter{enumi}{1}
\item
  拓扑群 G 上左一致结构与右一致结构的定义，只用到 G 中 e 的邻域滤子 B 的性质 \((\mathbf{GV}_{\mathbf{I}})\) 与 \((\mathbf{GV}_{\mathbf{II}})\) 。设 B 是满足这两条公理但不一定满足公理 \((\mathbf{GV}_{\mathbf{III}})\) 的一个滤子，试证左一致结构与右一致结构分别相容于 §1 习题 5 a) 中定义的拓扑 \(C_{s}\) 与 \(C_{d}\) 。
\item
  试证拓扑群 G 的下列条件等价：
\end{enumerate}

α) G 上的左一致结构与右一致结构相同。

\(\beta)\) 若 \(\mathbf{V}\) 是 \(e\) 在 \(\mathbf{G}\) 中的任一邻域，则存在一个邻域 \(\mathbf{W}\) （\(e\) 的），使得 \(x\mathbf{W}x^{-1} \subset \mathbf{V}\) 对一切 \(x \in \mathbf{G}\) 成立。

\(\gamma)\) 对称 \(x \to x^{-1}\) 是 \(\mathbf{G}_d\) 到 \(\mathbf{G}_d\) 上（或 \(\mathbf{G}_s\) 到 \(\mathbf{G}_s\) 上）的一致连续映射。

\(\delta)\) 映射 \((x,y)\to xy\) 是四个空间 \(\mathbf{G}_d\times \mathbf{G}_d,\mathbf{G}_s\times \mathbf{G}_s,\mathbf{G}_d\times \mathbf{G}_s,\mathbf{G}_s\times \mathbf{G}_d\) 之一到两个空间 \(\mathbf{G}_d,\mathbf{G}_s.\) 之一中的一致连续映射。

* 4) 设 \(\mathbf{G} = \mathbf{GL}(2, \mathbf{R})\) 是实元素的 \(2 \times 2\) 可逆方阵所成的乘法群。对每个整数 \(n > 0\) ，设 \(\mathbf{V}_n\) 是矩阵 \(X = \begin{pmatrix} x & y \\ z & t \end{pmatrix} \in \mathbf{G}\) （满足 \(|x - 1| \leqslant 1/n\) ，\(|y| \leqslant 1/n\) ，\(|z| \leqslant 1/n\) ，\(|t - 1| \leqslant 1/n\) ）组成的集合。试证诸 \(V_n\) 的族是单位元的一个邻域基本系（对应于一个与 \(G\) 的群结构相容的拓扑）。关于这个拓扑，\(G\) 是局部紧的，且 \(G\) 上的左一致结构与右一致结构是不同的。 *

\begin{enumerate}
\def\labelenumi{\arabic{enumi})}
\setcounter{enumi}{4}
\item
  设 G 是一个拓扑群，K 是 G 的一个正规子群。把 G/K 看作 \(\mathfrak{P}(G)\) 的一个子集，试证拓扑群 G/K 上的右一致结构是由 \(\mathfrak{P}(G)\) 上的一致结构诱导的，后者由 G 上的右一致结构按第二章 §1 习题 5 的程序构造。
\item
  试证拓扑群 G 上右一致结构与左一致结构的上确界（第二章，§ 2，第 5 节）是一个与 G 的拓扑相容的一致结构：称之为 G 上的双侧一致结构。试证每个豪斯多夫拓扑群同构于一个拓扑群的稠密子群，该拓扑群对于其双侧一致结构是完备的。
\end{enumerate}

¶ 7) a) 设 \(\mathbf{G}\) 是一个豪斯多夫拓扑群。若存在 \(\mathbf{V}_0\) （\(e\) 在 \(\mathbf{G}\) 中的一个邻域），使得对称 \(x \to x^{-1}\) （视为 \(\mathbf{G}_d\) 到 \(\mathbf{G}_d\) 中的映射）在 \(\mathbf{V}_0\) 上一致连续，试证 \(\mathbf{G}\) 有一个完备化（试证 § 4 定理 1 的条件满足）。

* b) 试证：若 G 是群的无限积，其中每个群都与习题 4 中定义的拓扑群相同，则 G 是完备的，但不存在 e 的邻域，使 \(x \rightarrow x^{-1}\) （视为 \(G_{d}\) 到 \(G_{d}\) 中的映射）在其上一致连续 *。

\begin{enumerate}
\def\labelenumi{\arabic{enumi})}
\setcounter{enumi}{7}
\item
  称豪斯多夫拓扑群 G 是局部准紧的，如果存在 e 的一个邻域 \(V_{0}\) ，它相对于 G 上的右（或左）一致结构是准紧的。试证每个局部准紧群都有一个局部紧的完备化（利用习题 7）。
\item
  设 G 是一个豪斯多夫拓扑群，K 是 G 的一个闭正规子群。试证：若拓扑群 K 与 G/K 都完备，则 G 也完备{[}考虑 G 上关于左（或右）一致结构的一个极小 Cauchy 滤子，以及它在 G/K 中的像{]}。
\item
  设 \((\mathbf{G}_i)_{i\in I}\) 是拓扑群的一个族，并设 \(\mathbf{G} = \prod_{i\in I}\mathbf{G}_i\) 是赋以 §2 习题 23 中定义的拓扑的积群。试证：若诸 \(\mathbf{G}_i\) 都完备，则 \(\mathbf{G}\) 也完备。
\item
  在 § 2 习题 26 的假设与记号下，设每个群 \(\mathbf{G}_{i}\) 都是豪斯多夫的且有完备化。试证局部积 \(\mathbf{G}\) （诸 \(\mathbf{G}_{i}\) 的，相对于开正规子群的一个族 \((\mathbf{K}_{i})\) ）有一个完备化 \(\hat{\mathbf{G}}\) ，它同构于诸 \(\hat{\mathbf{G}}_{i}\) 的局部积（相对于诸 \(\mathbf{K}_{i}\) 的闭包 \(\mathbf{K}_{i}\) ）。
\item
  \begin{enumerate}
  \def\labelenumii{\alph{enumii})}
  \tightlist
  \item
    设 \(G'\) 是一个完备豪斯多夫群；设 \(G_{0}\) 是 \(G'\) 的一个稠密子群，异于 \(G'\) ；并设 G 是给 \(G_{0}\) 赋以离散拓扑所得到的拓扑群。则恒等映射 \(G \rightarrow G_{0}\) 是一个双射连续同态，但它的连续扩张 \(\hat{G} \rightarrow \hat{G}_{0}\) 不是满射的。
  \end{enumerate}
\end{enumerate}

* b) 设 G 是群 \(\mathbf{Q}^2\) ，\(\theta\) 是一个无理数，\(u\) 是 G 到 R 中的连续同态 \((x,y)\to x + \theta y\) ，并设 \(\mathrm{G}' = u(\mathrm{G})\) 。则 \(u: \mathrm{G}\to \mathrm{G}'\) 是双射的，但它的连续扩张 \(\hat{\mathrm{G}}\to \hat{\mathrm{G}}'\) 不是单射的。

\begin{enumerate}
\def\labelenumi{\alph{enumi})}
\setcounter{enumi}{2}
\tightlist
\item
  由 a) 与 b) 构造一个双射连续同态 \(\mathbf{G} \to \mathbf{G}'\) （豪斯多夫交换群之间的）的例子，其连续扩张 \(\hat{\mathbf{G}} \to \hat{\mathbf{G}}'\) 既非单射也非满射*。
\end{enumerate}

